\documentclass[10pt, twocolumn]{article}
\usepackage[margin=0.75in]{geometry}
\usepackage{amsmath, amssymb}
\usepackage{graphicx}
\usepackage{booktabs}
\usepackage{hyperref}
\usepackage{abstract}

\renewcommand{\abstractnamefont}{\normalfont\bfseries}
\renewcommand{\abstracttextfont}{\normalfont\small\itshape}

\title{\textbf{REMORA \\ Real-time Expressive Motion to Output Routing Audio}}

\author{
    DRAVET Timothée$^{1}$, OVERHOLT Daniel$^{2}$ \\
    \small Department of Sound and Music Computing, Aalborg University \\
    \small $^{1}$\texttt{timothee.dravet@etu.univ-amu.fr} \\
    \small $^{2}$\texttt{dano@create.aau.dk}
}

\date{}

\begin{document}
\maketitle

\begin{onecolabstract}
This paper presents a two-column research paper template suitable for IEEE-style journal and conference submissions. The template demonstrates proper formatting for title, authors, abstract, sections, equations, tables, and references in a compact two-column layout commonly used in computer science and engineering publications.
\end{onecolabstract}

\noindent\textbf{Keywords:} two-column, research paper, LaTeX template, IEEE style

\section{Introduction}

Two-column layouts are standard in many conferences and journals, particularly in computer science, engineering, and physics. This template provides a starting point for such submissions.

The key advantages of two-column layouts include:
\begin{itemize}
    \item More efficient use of page space
    \item Better readability for technical content
    \item Conformance with venue requirements
\end{itemize}

\section{Related Work}

Discuss prior work relevant to your research. Cite important papers and explain how your approach differs or improves upon existing methods.

\section{Methodology}

\subsection{Problem Definition}

Given a dataset $\mathcal{D} = \{(\mathbf{x}_i, y_i)\}_{i=1}^{N}$, we aim to learn a function $f_\theta$ that minimizes:
\begin{equation}
    \mathcal{L}(\theta) = \frac{1}{N} \sum_{i=1}^{N} \ell(f_\theta(\mathbf{x}_i), y_i) + \lambda \|\theta\|_2^2
    \label{eq:loss}
\end{equation}

where $\ell$ is the loss function and $\lambda$ controls regularization.

\subsection{Proposed Approach}

Describe your method in detail. Use equations, algorithms, and diagrams as needed to explain the approach clearly.

\section{Experiments}

\subsection{Setup}

We evaluate on three benchmark datasets. All experiments use the same hyperparameters: learning rate $\eta = 10^{-3}$, batch size 32, and 100 training epochs.

\subsection{Results}

\begin{table}[h]
    \centering
    \caption{Comparison with baseline methods.}
    \label{tab:results}
    \begin{tabular}{lcc}
        \toprule
        Method & Acc. (\%) & F1 \\
        \midrule
        Baseline A & 82.1 & 0.79 \\
        Baseline B & 84.5 & 0.82 \\
        Baseline C & 86.3 & 0.84 \\
        \textbf{Ours} & \textbf{91.0} & \textbf{0.89} \\
        \bottomrule
    \end{tabular}
\end{table}

Table~\ref{tab:results} shows our method outperforms all baselines. The improvement is statistically significant ($p < 0.01$).

\section{Discussion}

Interpret the results and discuss implications, limitations, and potential future directions.

\section{Conclusion}

Summarize contributions and key findings. Our method achieves state-of-the-art performance on all evaluated benchmarks as shown in Equation~\ref{eq:loss} and Table~\ref{tab:results}.

\bibliographystyle{plain}
\bibliography{references}

\end{document}
